\documentclass[11pt]{article}
\usepackage[margin=25mm]{geometry}
\usepackage{fontspec}
\setmainfont{Latin Modern Roman}
\usepackage{amsmath,amssymb,mathtools}
\usepackage{unicode-math}
\setmathfont{Latin Modern Math}
\usepackage{xurl,hyperref}
\hypersetup{colorlinks=true,urlcolor=blue}
\setcounter{secnumdepth}{0}
\setlength{\parindent}{0pt}
\setlength{\parskip}{5pt}
\setlength{\emergencystretch}{3em}
\newcommand{\sref}[2]{\hyperlink{source:#1:#2}{[S#2]}}
\newcommand{\cataloguescope}[1]{\paragraph{Recorded target.}#1}
\newcommand{\R}{\mathbb R}
\newcommand{\C}{\mathbb C}
\newcommand{\N}{\mathbb N}
\newcommand{\Z}{\mathbb Z}
\newcommand{\Q}{\mathbb Q}
\newcommand{\F}{\mathbb F}
\newcommand{\eps}{\varepsilon}
\begin{document}
% UnsolvedMath ID 30000464; source catalogue code OWR-1266-002
\setcounter{section}{30000463}
\section{Nonperturbative Infrared Behavior of Nonamenable Sigma Models}\label{30000464}
{\small UnsolvedMath ID 30000464 · Mathematical Physics}
\subsection{Definitions and mathematical statement}

\paragraph{Shared notation.}$\mathbb N=\{1,2,\ldots\}$; the standard number fields have their usual meanings. All additional conventions are stated locally.


\paragraph{Definitions and conventions.} A nonlinear sigma model has fields taking values in a target manifold and an action determined by its target metric. Noncompact hyperbolic targets can have nonamenable symmetry groups, meaning that they lack an invariant mean on bounded functions. Infrared behavior concerns long spatial scales; an infrared-free perturbative RG prediction says that the effective coupling tends to zero at those scales. Nonperturbative analysis seeks conclusions about the full model rather than a formal finite-order series.

\paragraph{Source scope and qualifications.} There is a provenance mismatch: the supplied original excerpt is the functional-integral problem from Niedermaier's warped-product, dimensionally reduced gravity contribution, whereas the clean statement uses the infrared question in Seiler's separate nonamenable-target contribution. They must not be silently identified. The catalogue clean formulation is retained and both source targets are distinguished.

\paragraph{Original-excerpt target.} In the warped-product sigma models of the dimensionally reduced gravity discussion, develop the nonperturbative functional-integral formulation and relate it to the proposed algebraic bootstrap construction. This is separate from checking infrared freedom of the nonamenable-target model.

\paragraph{Catalogue mathematical target.}
\paragraph{Catalogue wording retained for traceability.}
For two-dimensional nonlinear sigma models with noncompact target spaces and nonamenable symmetry groups, does a rigorous nonperturbative analysis confirm the perturbative renormalization-group prediction of infrared asymptotic freedom, or is that prediction an artifact of perturbation theory?

\paragraph{Saved literature context.} The supplied catalogue's assessment dated 2026-08-21 records: The statement is a broad foundational question about nonperturbative functional integrals rather than a specific theorem; subsequent constructive results address special models but do not settle the general problem. The assessment is retained as dated provenance; it is not a new comprehensive literature or open-status certification.
\subsection{Short English statement}
Determine rigorously whether noncompact sigma models really have the long-distance behavior predicted by perturbation theory.
\subsection{Research attempt}
\paragraph{Model, date and investigated scope.}
OpenAI GPT-6 Astra (configured model \texttt{gpt-6-astra}), ultra reasoning; generated 11 October 2026. Constructs a positive-covariance beyond-all-orders correction that changes long-distance clustering, testing the perturbative inference without asserting a sigma-model example.

\paragraph{Sources and literature check.}
The complete supplied definition, including source qualifications and the dated literature assessment, was reviewed. The calculation below is an elementary restricted attempt motivated by [S2]; it does not claim novelty or a fresh literature-status determination.

\paragraph{Context and investigated scope.} The complete supplied definition, including source qualifications and the dated literature assessment, was reviewed. The calculation below is an elementary restricted attempt motivated by [S2]; it does not claim novelty or a fresh literature-status determination.

\paragraph{Testing what a perturbation series can decide.} Let $0<g<1$ and consider an arbitrary stationary centered Gaussian field with covariance $C_0(x-y)$ tending to zero at large separation. Add an independent, spatially constant centered Gaussian variable of variance $e^{-1/g}$. The resulting covariance is
\[
 C_g(x-y)=C_0(x-y)+e^{-1/g}.
\]
It is positive definite: for coefficients $a_j$ the additional quadratic form is $e^{-1/g}|\sum_ja_j|^2\geq0$. All derivatives of $e^{-1/g}$ at $g=0^+$ vanish, since they are that exponential times a polynomial in $1/g$. Consequently its formal expansion agrees to every order with that of $C_0$, but $\lim_{|x-y|\to\infty}C_g(x-y)=e^{-1/g}>0$ for every fixed $g>0$.

\paragraph{Nonuniformity is the obstruction.} If $C_0(r)=r^{-\eta}$ in a range with $\eta>0$, the added term dominates beyond $r>e^{1/(\eta g)}$. No finite-order or even full formal expansion at fixed distance rules this out. This is an explicitly constructed covariance countertest to a method, not a nonamenable-target sigma-model counterexample.

\paragraph{Both source targets remain distinct.} To establish or refute infrared freedom in the actual two-dimensional noncompact model, one needs estimates uniform on such coupling-dependent distance scales and control of the noncompact zero mode. The separate warped-product functional-integral/ bootstrap question in the original excerpt additionally requires construction of a measure and comparison of its observables with algebraic data. Neither follows from the covariance example, and the two provenance targets have not been conflated. The first gap in the infrared route is uniform nonperturbative control, not a further coefficient of the perturbation series.

\paragraph{Verification and adversarial audit.}
Checked covariance positivity, flatness at zero coupling, noncommuting limits, and the two different source targets. Model self-audit only; no independent human proof checking. Prompts A and B were applied to this canonical record and its qualified source scope. Checked covariance positivity, flatness at zero coupling, noncommuting limits, and the two different source targets. Retained unresolved status and zero original-target progress; Uniform long-distance nonperturbative estimates are absent; the separate warped-product measure/bootstrap construction also remains unresolved.

Coordinator checked the positive constant Gaussian covariance perturbation, its flat dependence at zero coupling and its nondecaying large-distance plateau. The constructed covariance is a diagnostic example. This additional check was model-assisted, not independent human verification.

\paragraph{Tools, process and contributions.}
Python 3 was used for file preparation and structural checks; no numerical experiment or theorem prover was used. Contribution prompt A was followed by prompt B. GPT-6 Astra with ultra reasoning drafted the research and performed its stated self-audit; any additional coordinator contribution is identified in the verification and audit. A GPT-6-based Codex coordinator assembled the manuscript and checked repository integration; its exact variant and reasoning configuration were not exposed. No human mathematical contribution, independent human proof check or formal proof certification is supplied. Complete derivations and the consulted source references are included here; private reasoning transcripts are not published.

\paragraph{Final status and first remaining gap.}
\textbf{UNRESOLVED}. Uniform long-distance nonperturbative estimates are absent; the separate warped-product measure/bootstrap construction also remains unresolved.

\section{Completion Estimate}
\noindent\textbf{COMPLETION: 0\%}
The calculation is an auxiliary or restricted attempt and establishes no newly verified portion of the full source target. Uniform long-distance nonperturbative estimates are absent; the separate warped-product measure/bootstrap construction also remains unresolved.

\subsection{Sources}
\begin{itemize}
\item[\hypertarget{source:30000464:1}{S1}] ULAM.AI and the UnsolvedMath Contributors, UnsolvedMath, record 30000464 (OWR-1266-002). Catalogue curation under CC BY 4.0; definitions and English exposition added locally.
\url{https://huggingface.co/datasets/ulamai/UnsolvedMath}.
\item[\hypertarget{source:30000464:2}{S2}] The Rigorous Renormalization Group (2006), Oberwolfach Reports, DOI 10.4171/owr/2006/17, pp. 1027--1076; official report source and mathematical context.
\url{https://doi.org/10.4171/owr/2006/17}.
\item[\hypertarget{source:30000464:3}{S3}] Official EMS Press PDF of the same Oberwolfach report; PDF pages 15 and 42--43, printed pages 1041 and 1068--1069.
\url{https://ems.press/content/serial-article-files/46049}.
\end{itemize}
\label{end:30000464}
\end{document}
